\chapter{技術士法・義務責務}
技術士法の第4章は、専門家が高度な知識を持つことに対応した行動上の責任を定める。数式の正しさだけで技術サービスの適否は決まらない。技術の危険を十分に評価し、関係者へ事実を説明し、資格と能力の範囲を明らかにすることが信頼の前提である。本章では条文の語句、具体的な業務判断、資格取得後の継続研鑚を結び付ける。

\section{義務・責務を主語と条件から読む}
\label{theory:engineer_act:duties}
過去問で示された第4章の条文と、2026年10月8日にe-Gov公式APIから取得した2025年6月1日施行版を対照した。以下は本章の設問で扱う範囲の整理であり、全改正履歴の照合を意味しない\cite{engineer_act_current}。
\begin{longtable}{p{.13\linewidth}p{.29\linewidth}p{.47\linewidth}}
条文 & 要点 & 適用対象と条件\\\hline
第44条 & 信用失墜行為の禁止 & 技術士・技術士補。専門家集団の信用を害する行為を禁止する。\\
第45条 & 秘密保持義務 & 業務上知った秘密の漏えい・盗用を正当な理由なしに行わない。資格を失った後も続く。\\
第45条の2 & 公益確保の責務 & 技術士・技術士補。公共の安全、環境その他公益を害さないよう努める。\\
第46条 & 名称表示の場合の義務 & 技術士の名称を業務に関して表示する場合、登録技術部門を明示し、未登録部門を表示しない。\\
第47条 & 技術士補の業務の制限等 & 技術士補の名称での業務は技術士を補助する場合に限る。名称表示の規定を準用する。\\
第47条の2 & 資質向上の責務 & 技術士。その業務の知識技能と資質を継続的に高めるよう努める。\\
\end{longtable}
第4章という見出し全体には「義務」と書かれるが、個々の条文には義務と責務の両方がある。公益確保・資質向上は、努力を求める責務という語で整理する。「責務だからしなくてもよい」という結論は出ない。名称表示をする場合という条件を、すべての業務実施時に無条件に拡張することも避ける。

法令で区分された登録先は「技術部門」である。「専門部門」という日常語へ置き換えない。第47条2項の「準用」は、別の対象へ必要な読み替えを伴って規定を働かせることであり、その規定が元から直接対象とする場合の適用とは区別する。平成26・29・30年度、令和元・2・6年度の空欄問題を、この表と照合すると選択肢の違いを説明できる。

登録部門の力量は取得時に固定されるものではない。技術の変化、法令、失敗事例を学び、自分の業務能力を更新する必要がある。別部門を学ぶことは有益だが、現在の登録部門の学習を不要にするものではない。また顧客の同意だけでは技術士補の名称による独立業務や未登録技術部門の表示が認められることにはならない。

\section{守秘・公益・説明責任を同時に満たす}
\label{theory:engineer_act:public}
守秘義務は、情報を託す相手が安心して技術者に相談できるための基盤である。所属組織の取引先だけでなく、業務上知った秘密一般を対象とし、退職・資格喪失後も存続する。一方、公衆に危険がある場合にも顧客の利益だけを優先することはできない。技術士倫理綱領は、公衆の利益、有能性、真実性、秘密の保持などを相互に関連する行動原則として挙げる\cite{cpd_v3}。

例えば試験データに危険な不具合が表れているのに顧客が報告から除外するよう指示した場合、まず事実と不確かさを整理し、危険と必要な対策を責任ある相手へ説明する。組織内で解決しないときも、公衆の安全を確保する適切な手段を検討する。「顧客命令なら何でも従う」「守秘を理由に危険を放置する」「すべての秘密を直ちに公開する」という一律の処理では解けない。公益確保と守秘の対象・理由・手続を具体的に検討する必要がある。

利益相反は、報酬・所属・人的関係などが専門的判断の公正さに影響し得る状態である。関係を説明し、判断の根拠を客観的な事実へ戻す。関係者の話をよく聞くことと、事実に反した結論を受け入れることは別である。

依頼者の情報から独自に調査したとしても、元の秘密性や契約上の利用制限が自動的に消えるわけではない。さらに情報の所有・利用権と守秘義務は別の問題である。平成23年度問1では権利を当事者の協議へ委ねる表現、平成27年度問1・平成30年度問2・令和5年度問1では独自調査だから制限がなくなるという表現の違いを読む。

\section{適性科目の規範：法の形式と倫理の内容}
\label{theory:engineer_act:norms}
社会規範は、人々の行動を社会の中で評価する基準である。そのうち法規範は法律として制定され、倫理規範は何が責任ある行動かを示す。技術士法第4章は形式上の法律であり、同時に専門家の職業倫理を表している。この二つの性格は両立する。令和元年度再試験問1は、社会規範を上位概念とし、第4章を法規範かつ倫理規範として整理する問題である。

\section{CPD：資格取得後に能力を更新する}
\label{theory:engineer_act:cpd}
CPDは Continuing Professional Development の略である。学習・研修・調査研究などを通して専門家としての能力を継続して高める。指定された第2版は「自己研鑚」、第3版は「継続研鑚」という語を用いる。両版の2.2節では、その対象が本人の能力と立場によって異なり、実施記録と証拠により説明できることを求めている\cite{cpd_v2,cpd_v3}。平成27年度問2・平成30年度問3の空欄は、出題指定資料の語句を確かめる。

現在の知識を使って日常業務を繰り返すことと、その業務に関連して新しい方法を調べ能力を向上させることは別である。研修で何を学んだか、業務でどの能力が増えたかを記録する。実施の事実があっても、その内容を証明できなければ登録上の実績として認められない場合がある。

2021年の大臣通知・省令改正に基づく制度では、希望する技術士のCPD実績を登録簿の資質向上の取組状況へ記載し、条件に応じた認定や公表を行う仕組みが整えられた\cite{cpd_new2021}。令和6年度問15はこの制度も問う。実績の記載・認定制度を、すべての技術士に一律の登録更新義務があるという制度へ読み替えない。

ガイドラインと算定基準は実施年度によって異なる。第1版の対象は2011〜2013年度、第2版は2014〜2016年度、第3版は2017〜2021年度であり、その後の新制度と現在の版は別である\cite{cpd_versions}。過去問を現在の時間算定基準だけで解かない。

平成23年度問4では、資質向上の不履行を「行政罰」に相当するとした表現が複数正答の検討点になる。第36条2項の登録取消し等という行政上の措置、罰則、努力責務の履行評価を区別し、原典の記述と公式の複数正答を保持する。どのようなCPD不足でも自動的に登録が取り消されるという説明にはしない。

\section{コンピテンシー：知識を仕事の能力へ結び付ける}
\label{theory:engineer_act:competency}
平成26年3月の技術士分科会資料は、専門的学識、問題解決、マネジメント、評価、コミュニケーション、リーダーシップ、技術者倫理の7つを示す\cite{mext_competency2014}。令和元年度問2では、列挙された7つすべてがその資料に含まれる。

物理の専門知識を正確に適用するだけでは、複数の制約がある業務を完遂できない。問題解決では矛盾する要求を整理し、マネジメントでは人・設備・資金・情報を配分し、評価では成果を次の改善へつなげる。コミュニケーションとリーダーシップは関係者の理解と調整、技術者倫理は公益と責任を支える。CPDはこれらの能力を業務の変化に合わせて更新する活動である。ここでは出題指定の平成26年資料の項目を扱い、その後の全改訂を同一と見なさない。

\section{教育認定から国際的な業務資格へ}
\label{theory:engineer_act:international}
技術者が国境を越えて活動するには、教育の質、実務能力、資格登録、受入国の法制度を区別する必要がある。教育課程の認定は、その課程で得られる学習成果の質を評価する。JABEEの認定だけで全ての受入国の業務免許が与えられるわけではない。

\begin{description}
 \item[教育と第一次試験免除] 文部科学大臣が指定する認定課程の修了者は、技術士補となる資格を持つ。第一次試験を免除されることと、技術士補として登録を受けることは別の段階である\parencite{engineer_act_current,ipej_associate_registration}。
 \item[国際称号] IPEAやAPECの登録では、教育・実務経験・独立した業務遂行能力・継続学習等の要件を審査する。国内の技術士資格だけで自動的に国際登録となるのではない。IPEAの和文名称は2015年4月に旧EMF国際エンジニアから変更された\parencite{ipej_ipea2014}。
 \item[受入国の業務資格] 共通称号は能力の同等性を示す枠組みである。相互承認には別の二国間枠組みや固有事項審査等が関わり、受入国の法・登録・入国手続まで全て不要になるとはいえない。
\end{description}

\paragraph{平成27年当時と現在の建築教育}
建築系の認定についてJABEEは2012年のUNESCO-UIAシステム認証、2014年のキャンベラ協定暫定加盟、2019年の正式加盟を説明する\parencite{jabee_canberra}。平成27年の肢にあるUNESCO-UIAの要約を現在の正式加盟先に置き換えず、機関の認証と協定への加盟という手続の違いも確認する。

\paragraph{二国間の審査合理化}
文科省は日豪の2003年枠組み、建築技術教育普及センターは2008年の日豪・2009年の日NZ枠組みを説明する\parencite{mext_international,jaeic_mutual2010}。一般的な能力の審査と、受入国固有の法的・技術的事項の審査を分ける。現在の手続や対象分野の表を、2015年の全条件と同一だと扱わない。

\section*{各問題の個別検討}
以下は査読待ち初稿である。平成25年度問1では独立した条文読解と公式の個数が一致しないため、相違を表示する。平成23年度問4の複数正答も一つに確定しない。
\section{継続研鑽：活動・実績・証明をつなぐ}
\label{theory:profession:cpd}
CPDは専門職として能力を継続的に維持・向上させる活動である。仕事に従事している事実だけから、必要な学習と実績記録をすべて満たすとはしない。実務の中で何を学び、能力をどう改善したかを説明できる記録を残す\parencite{cpd_v3,cpd_versions,cpd_new2021}。

令和4年度問題には参加型、発信型、実務型、自己学習型という区分がある。研修の受講、技術指導や執筆、業務成果、自己学習は活動の性格が違う。換算条件と証明は制度・版ごとに確認する。学協会によって資格継続に確認条件があることと、国内の技術士資格が一律に自動失効することも区別する。

\section{2023コンピテンシーの八つのキーワード}
\label{theory:profession:competency2023}
2023年の資料では、専門的学識、問題解決、マネジメント、評価、コミュニケーション、リーダーシップ、技術者倫理、継続研さんの八つを挙げる。2014年の七つのキーワードから継続研さんを追加した時点を確認する\parencite{eth54_competency2023}。

知識を持つことだけでなく、複合的な問題を定義し、制約と利害を調べ、資源を配分し、結果を評価し、関係者と協力して責任ある成果を得る能力を考える。35歳程度での活躍への期待は、法定の受験年齢上限や資格の自動取消条件ではない。文書の版と発表主体を読み、2023版で独立しているコミュニケーションとリーダーシップを一つの項目にまとめない。

\subsection{平成23年度 問1}
\label{q:h23-reg-apt-01}

\noindent\textbf{原典：} sources/H23/aptitude.pdf、PDF物理2頁。

\textbf{公式正答：} 4（single）。

\par\textit{原典本文の校正と独自解説は未作成。分類のみの登録項目。}

\subsection{平成23年度 問4}
\label{q:h23-reg-apt-04}

\noindent\textbf{原典：} sources/H23/aptitude.pdf、PDF物理5頁。

\textbf{公式正答：} 1|5（multiple\_accepted）。

\par\textit{原典本文の校正と独自解説は未作成。分類のみの登録項目。}

\subsection{平成24年度 問1}
\label{q:h24-reg-apt-01}

\noindent\textbf{原典：} sources/H24/aptitude.pdf、PDF物理1頁。

\textbf{公式正答：} 5（single）。

\par\textit{原典本文の校正と独自解説は未作成。分類のみの登録項目。}

\subsection{平成24年度 問14}
\label{q:h24-reg-apt-14}

\noindent\textbf{原典：} sources/H24/aptitude.pdf、PDF物理12頁。

\textbf{公式正答：} 4（single）。

\par\textit{原典本文の校正と独自解説は未作成。分類のみの登録項目。}

\subsection{平成25年度 問1}
\label{q:h25-reg-apt-01}

\noindent\textbf{原典：} sources/H25/aptitude.pdf、PDF物理1頁。

\textbf{公式正答：} 4（single）。

\par\textit{原典本文の校正と独自解説は未作成。分類のみの登録項目。}

\subsection{平成26年度 問1}
\label{q:h26-reg-apt-01}

\noindent\textbf{原典：} sources/H26/aptitude.pdf、PDF物理1頁。

\textbf{公式正答：} 3（single）。

\par\textit{原典本文の校正と独自解説は未作成。分類のみの登録項目。}

\subsection{平成27年度 問1}
\label{q:h27-reg-apt-01}

\noindent\textbf{原典：} sources/H27/aptitude.pdf、PDF物理1頁。

\textbf{公式正答：} 5（single）。

\par\textit{原典本文の校正と独自解説は未作成。分類のみの登録項目。}

\subsection{平成27年度 問2}
\label{q:h27-reg-apt-02}

\noindent\textbf{原典：} sources/H27/aptitude.pdf、PDF物理2頁。

\textbf{公式正答：} 5（single）。

\par\textit{原典本文の校正と独自解説は未作成。分類のみの登録項目。}

\subsection{平成28年度 問1}
\label{q:h28-reg-apt-01}

\noindent\textbf{原典：} sources/H28/aptitude.pdf、PDF物理1頁。

\textbf{公式正答：} 2（single）。

\par\textit{原典本文の校正と独自解説は未作成。分類のみの登録項目。}

\subsection{平成28年度 問2}
\label{q:h28-reg-apt-02}

\noindent\textbf{原典：} sources/H28/aptitude.pdf、PDF物理2頁。

\textbf{公式正答：} 3（single）。

\par\textit{原典本文の校正と独自解説は未作成。分類のみの登録項目。}

\subsection{平成29年度 問1}
\label{q:h29-reg-apt-01}

\noindent\textbf{原典：} sources/H29/aptitude.pdf、PDF物理1頁。

\textbf{公式正答：} 3（single）。

\par\textit{原典本文の校正と独自解説は未作成。分類のみの登録項目。}

\subsection{平成29年度 問2}
\label{q:h29-reg-apt-02}

\noindent\textbf{原典：} sources/H29/aptitude.pdf、PDF物理2頁。

\textbf{公式正答：} 5（single）。

\par\textit{原典本文の校正と独自解説は未作成。分類のみの登録項目。}

\subsection{平成30年度 問1}
\label{q:h30-reg-apt-01}

\noindent\textbf{原典：} sources/H30/aptitude.pdf、PDF物理1頁。

\textbf{公式正答：} 2（single）。

\par\textit{原典本文の校正と独自解説は未作成。分類のみの登録項目。}

\subsection{平成30年度 問2}
\label{q:h30-reg-apt-02}

\noindent\textbf{原典：} sources/H30/aptitude.pdf、PDF物理2頁。

\textbf{公式正答：} 5（single）。

\par\textit{原典本文の校正と独自解説は未作成。分類のみの登録項目。}

\subsection{平成30年度 問3}
\label{q:h30-reg-apt-03}

\noindent\textbf{原典：} sources/H30/aptitude.pdf、PDF物理3頁。

\textbf{公式正答：} 3（single）。

\par\textit{原典本文の校正と独自解説は未作成。分類のみの登録項目。}

\subsection{令和元年度 問1}
\label{q:r01-reg-apt-01}

\noindent\textbf{原典：} sources/R01/aptitude.pdf、PDF物理1頁。

\textbf{公式正答：} 5（single）。

\par\textit{原典本文の校正と独自解説は未作成。分類のみの登録項目。}

\subsection{令和元年度 問2}
\label{q:r01-reg-apt-02}

\noindent\textbf{原典：} sources/R01/aptitude.pdf、PDF物理3頁。

\textbf{公式正答：} 5（single）。

\par\textit{原典本文の校正と独自解説は未作成。分類のみの登録項目。}

\subsection{令和元年度再試験 問1}
\label{q:r01-ret-apt-01}

\noindent\textbf{原典：} sources/R01\_retest/aptitude.pdf、PDF物理1頁。

\textbf{公式正答：} 3（single）。

\par\textit{原典本文の校正と独自解説は未作成。分類のみの登録項目。}

\subsection{令和元年度再試験 問2}
\label{q:r01-ret-apt-02}

\noindent\textbf{原典：} sources/R01\_retest/aptitude.pdf、PDF物理2頁。

\textbf{公式正答：} 3（single）。

\par\textit{原典本文の校正と独自解説は未作成。分類のみの登録項目。}

\subsection{令和2年度 問1}
\label{q:r02-reg-apt-01}

\noindent\textbf{原典：} sources/R02/aptitude.pdf、PDF物理1頁。

\textbf{公式正答：} 5（single）。

\par\textit{原典本文の校正と独自解説は未作成。分類のみの登録項目。}

\subsection{令和3年度 問1}
\label{q:r03-reg-apt-01}

\noindent\textbf{原典：} sources/R03/aptitude.pdf、PDF物理1頁。

\textbf{公式正答：} 3（single）。

\par\textit{原典本文の校正と独自解説は未作成。分類のみの登録項目。}

\subsection{令和4年度 問1}
\label{q:r04-reg-apt-01}

\noindent\textbf{原典：} sources/R04/aptitude.pdf、PDF物理1頁。

\textbf{公式正答：} 4（single）。

\par\textit{原典本文の校正と独自解説は未作成。分類のみの登録項目。}

\subsection{令和5年度 問1}
\label{q:r05-reg-apt-01}

\noindent\textbf{原典：} sources/R05/aptitude.pdf、PDF物理1頁。

\textbf{公式正答：} 3（single）。

\par\textit{原典本文の校正と独自解説は未作成。分類のみの登録項目。}

\subsection{令和6年度 問1}
\label{q:r06-reg-apt-01}

\noindent\textbf{原典：} sources/R06/aptitude.pdf、PDF物理1頁。

\textbf{公式正答：} 4（single）。

\par\textit{原典本文の校正と独自解説は未作成。分類のみの登録項目。}

\subsection{令和6年度 問15}
\label{q:r06-reg-apt-15}

\noindent\textbf{原典：} sources/R06/aptitude.pdf、PDF物理15頁。

\textbf{公式正答：} 2（single）。

\par\textit{原典本文の校正と独自解説は未作成。分類のみの登録項目。}

\subsection{令和7年度 問1}
\label{q:r07-reg-apt-01}

\noindent\textbf{原典：} sources/R07/aptitude.pdf、PDF物理1頁。

\textbf{公式正答：} 3（single）。

\par\textit{原典本文の校正と独自解説は未作成。分類のみの登録項目。}

